\chapter*{درباره‌ی این جزوه}
\addcontentsline{toc}{chapter}{درباره‌ی این جزوه}
\markboth{درباره‌ی این جزوه}{درباره‌ی این جزوه}

این جزوه بر پایه‌ی اسلایدهای درس \textbf{جبر خطی} (\en{CE282}) دانشکده‌ی مهندسی کامپیوتر دانشگاه صنعتی شریف نوشته شده و مطالب را به همان ترتیب اسلایدها، با زبان ساده توضیح می‌دهد. هر جا در اسلایدها جای اثبات خالی مانده بود، اثبات نوشته شده و مثال‌های بی‌پاسخ حل شده‌اند. چند خطای کوچک اسلایدها هم اصلاح شده و در پانویس آمده است.

\section*{راهنمای جعبه‌ها}
\begingroup
\renewcommand{\theboxnum}{۲-۵}
\begin{definition}[نام مفهوم]
تعریف یک مفهوم جدید، همراه با واژه‌ی انگلیسی آن.
\end{definition}
\begin{theorem}[نام قضیه]
یک حکم مهم؛ اثباتش در جعبه‌ی خاکستری بعدی می‌آید.
\end{theorem}
\begin{example}
مثال‌ها نارنجی‌اند و پاسخشان زیرشان می‌آید.
\end{example}
\endgroup
بقیه‌ی جعبه‌ها: \textcolor{lemcolor}{\semifont لم و نتیجه}، \textcolor{notecolor}{\semifont نکته}، \textcolor{warncolor}{\semifont مراقب باشید!} (اشتباه‌های رایج)، \textcolor{intuitcolor}{\semifont شهود} و \textcolor{sumcolor}{\semifont جمع‌بندی} (پایان هر فصل).

\section*{نمادگذاری}
\begin{itemize}
  \item بردارها پررنگ‌اند ($\vx,\vb$) و به‌طور پیش‌فرض ستونی؛ $\zero$ بردار صفر است.
  \item ماتریس‌ها با حروف بزرگ ($A$) و درایه‌ها با $A_{ij}$ یا $a_{ij}$ نوشته می‌شوند؛ $A\T$ ترانهاده است.
  \item «قضیه‌ی ۲-۵» یعنی جعبه‌ی پنجم فصل ۲.
\end{itemize}
